\begin{pastexam}[label={px:ncku-114-ee-ds-12}]{成大電機所 114 學年度 單選題第 12 題}
（單選；每題 5 分，答錯倒扣 2 分）

Continue from the problem 11, instead of chaining, we adopt open addressing to handle overflow,
where each bucket has one slot. The resulting position of data is determined by h(key)=
\{initial position + $i$* h$_2$(key) \}\%11, where h$_2$(key) = key \% 7 and $i$ denotes the
iteration of the overflow. Please show the position of 58 in the table.

a) 7 \quad b) 5 \quad c) 10 \quad d) 3 \quad e) None of the above.
\end{pastexam}

\begin{pastsoln}
答案是 a) 7。

這是雙重雜湊（資結書第~12.3~節）：初始位置 $h_1 = \mathit{key} \bmod 11$，第 $i$ 次碰撞後試
$(h_1 + i \cdot h_2) \bmod 11$，$h_2 = \mathit{key} \bmod 7$。 照第 11 題的順序插入，到 $58$ 為止：
\begin{center}
\begin{tabular}{cccll}
\toprule
鍵值 & $h_1$ & $h_2$ & 試的位置 & 放在 \\
\midrule
3 & 3 & & 3 & 3 \\
41 & 8 & & 8 & 8 \\
15 & 4 & & 4 & 4 \\
36 & 3 & $36 \bmod 7 = 1$ & 3（滿），$3 + 1 = 4$（滿），$3 + 2 = 5$ & 5 \\
74 & 8 & $74 \bmod 7 = 4$ & 8（滿），$(8 + 4) \bmod 11 = 1$ & 1 \\
58 & 3 & $58 \bmod 7 = 2$ & 3（滿），$3 + 2 = 5$（滿），$3 + 4 = 7$ & \textbf{7} \\
\bottomrule
\end{tabular}
\end{center}
$58$ 放在第 $7$ 格。

原題的雜湊設計有個問題：$h_2(\mathit{key}) = \mathit{key} \bmod 7$ 可能是 $0$。 下一個要插入的
$91 = 7 \times 13$，$h_2 = 0$，而 $h_1 = 3$ 已經被佔，之後每一次都還是試第 $3$ 格，永遠放不進去。
雙重雜湊的第二個雜湊函數不能是 $0$，而且要和表的大小互質，常見的寫法是
$h_2(\mathit{key}) = 1 + (\mathit{key} \bmod 7)$ 這類形式。 這不影響本題問的 $58$，因為它比 $91$
早插入。
\end{pastsoln}
